\documentclass[11pt]{article}

\usepackage[a4paper,margin=1in]{geometry}
\usepackage{amsmath, amssymb, amsthm}
\usepackage{microtype}
\usepackage{newpxtext, newpxmath}
\usepackage{xcolor}
\usepackage[most]{tcolorbox}
\usepackage{enumitem}
\usepackage{titlesec}

\definecolor{inkblue}{HTML}{1F3A5F}
\definecolor{boxbg}{HTML}{F2F5F9}
\definecolor{boxrule}{HTML}{8AA4C4}

\titleformat{\section}{\large\bfseries\color{inkblue}}{}{0em}{}
\titlespacing{\section}{0pt}{1.4em}{0.5em}

\newtcolorbox{insight}{
  colback=boxbg, colframe=boxrule, boxrule=0.6pt,
  arc=2pt, left=10pt, right=10pt, top=8pt, bottom=8pt,
  breakable, enhanced,
  title={\textbf{Insight}}, coltitle=inkblue, fonttitle=\bfseries\small,
  attach title to upper=\hspace{0.6em}
}

\newtcolorbox{summary}{colback=yellow!18, colframe=yellow!55!black, arc=2pt, boxrule=0.5pt,
  breakable, left=10pt, right=10pt, top=8pt, bottom=8pt}

\setlength{\parskip}{0.55em}
\setlength{\parindent}{0pt}

\title{\color{inkblue}\bfseries Surprise}
\author{}
\date{}

\begin{document}
\maketitle
\thispagestyle{empty}

\begin{summary}
A language model's loss is the average surprise it feels at the real next word,
and that average splits cleanly into two pieces: the entropy of language, which nothing can
remove, and the model's own avoidable mistakes --- which is the KL-divergence. So the loss and
the KL are the same quantity up to a fixed constant, and a single line of algebra shows why two
formulas that look nothing alike are the same thing.
\end{summary}

\section{Information: the surprise of a single outcome}

A language model, at each step, is placing bets. Before it sees the next word it spreads its
confidence across all the words it knows --- a little probability here, a lot there. Then the
real next word arrives, and we ask one question: \emph{how badly did that catch us off guard?}

If the outcome was one we had thought likely, we are barely surprised; if we had thought it
nearly impossible, we are very surprised. So we want a number --- call it the \textbf{surprise},
or the \emph{information content}, of an outcome --- that depends only on the probability $p$ we
had assigned it, and we want it to behave sensibly. Three demands pin it down completely.
\begin{enumerate}[leftmargin=1.6em]
  \item \textbf{Rarer is more surprising.} Surprise should \emph{decrease} as $p$ grows: a likely
  outcome shocks us less than an unlikely one.
  \item \textbf{Certainty is no surprise.} If $p = 1$ the outcome was guaranteed, and learning it
  tells us nothing, so its surprise should be exactly $0$.
  \item \textbf{Independent surprises add.} If two \emph{independent} things happen --- the first
  with probability $p$, the second with probability $q$ --- then learning both should surprise us
  by the \emph{sum} of the two separate surprises. Independent events carry separate,
  non-overlapping information, so their informations should simply pile up.
\end{enumerate}

The third demand is the one with teeth, and it is the whole reason a logarithm appears. For
independent events the probabilities \emph{multiply}: the chance of both is $p \cdot q$. Demand~3
says the surprise of that combined event, $S(p\,q)$, must equal $S(p) + S(q)$. So the surprise
function has to turn \emph{multiplication into addition}:
\[
  S(p \cdot q) \;=\; S(p) + S(q).
\]
There is essentially only one function that does this, and it is the logarithm --- that is
precisely the property $\log(pq) = \log p + \log q$. (One can see the shape forced directly:
putting $q = p$ gives $S(p^2) = 2\,S(p)$, and repeating, $S(p^n) = n\,S(p)$ --- surprise scales
with the \emph{exponent}, and reading off exponents is exactly what a logarithm does.) Combining
this with demand~1, that surprise decreases in $p$, fixes the sign: we need a \emph{minus} sign,
because $\log p$ itself \emph{increases} toward $0$ as $p$ grows to $1$. The result is the
definition:
\[
  \boxed{\,I(x) \;=\; -\log p(x)\,}
\]
the \textbf{information content} (or \emph{self-information}) of an outcome $x$ of probability
$p(x)$. Demand~2 comes for free: $I = -\log 1 = 0$ for a certain outcome. And as $p$ slides
toward $0$, $-\log p$ climbs without bound --- an outcome you were sure would not happen,
happening anyway, is an unbounded surprise, exactly as it should feel.

\begin{insight}
The logarithm is not a modelling choice; it is the \emph{only} way to meet demand~3. Surprise
lives in an additive world (informations pile up) while probability lives in a multiplicative one
(independent chances multiply), and the logarithm is the unique bridge between them. Keep this in
view: the \emph{same} trick returns later in this note, when $\log$ turns the ratio $p/q$ into the
difference $\log p - \log q$ and splits the KL-divergence in two. One function, one job, used
twice.
\end{insight}

\subsection*{Units: bits and nats}

A logarithm needs a base, and the base is just a choice of \emph{unit} for surprise. Base $2$
measures information in \textbf{bits}. A fair coin has $p = \tfrac12$ for each face, so each toss
carries
\[
  I = -\log_2 \tfrac12 = 1 \text{ bit},
\]
the textbook ``one bit'' --- the information in a single yes/no answer with no prior lean. A
\emph{biased} coin, say $p(\text{heads}) = 0.9$, is different: heads now costs only
$-\log_2 0.9 \approx 0.15$ bits (barely a surprise), while the rare tails costs
$-\log_2 0.1 \approx 3.32$ bits. Machine-learning work usually uses the natural logarithm
instead, measuring in \textbf{nats}; the two differ only by a fixed scale, $1 \text{ nat} =
1/\!\ln 2 \approx 1.443$ bits. Nothing below depends on the choice, so we write just $\log$ and
keep whichever unit is convenient.

\section{Entropy: the average surprise}

Information content is the surprise of \emph{one} outcome, known after the fact. But before the
outcome we can already ask: \emph{how surprised should we expect to be?} That expected surprise,
averaged over all outcomes by how often each occurs, is the \textbf{entropy} of the source:
\[
  H(X) \;=\; \mathbb{E}\bigl[I(X)\bigr] \;=\; -\sum_x p(x)\,\log p(x).
\]
It is the average information per observation --- the uncertainty built into $X$ before you look.

The fair coin has $H = 1$ bit: both outcomes carry one bit, so the average is one bit, the most
an un-tilted two-way choice can hold. The biased $0.9$ coin has
\[
  H = -\bigl(0.9\log_2 0.9 + 0.1\log_2 0.1\bigr) \approx 0.47 \text{ bits},
\]
less than one, because the coin is \emph{predictable}: you can usually guess heads, so you are
rarely surprised, and your average surprise is low. High entropy means genuinely uncertain
outcomes (the fair coin); low entropy means predictable ones (the biased coin). Entropy is the
average surprise you cannot talk your way out of --- it belongs to the source itself.

\begin{insight}
This is the object that becomes the ``built-in unpredictability of language'' in the sections
below. When we say no model, however good, can get beneath a certain floor of surprise, that
floor \emph{is} the entropy $H$ of the data: the average surprise that remains even when your
probabilities are perfectly correct. Entropy is the irreducible part; everything above it is the
model's own fault.
\end{insight}

\section{A model's loss is average surprise --- at the real word}

Now the model. For a context, it outputs a probability $q(w)$ for every possible next word $w$.
The true next word $y$ arrives, and the model's surprise is its own information content for that
word,
\[
  L \;=\; -\log q(y),
\]
zero if the model was certain and right ($q(y) = 1$), large if it had given the true word little
chance. This is the \emph{same} $-\log$ of surprise as before --- only now the probability inside
it is the \emph{model's} guess $q$, not the true probability.

Training runs the model over an enormous amount of real text and averages this surprise. Each
possible word $w$ shows up in real language some fraction of the time --- call that true frequency
$p(w)$ --- so the average surprise weights each word's cost $-\log q(w)$ by how often it really
occurs:
\[
  \text{loss} \;=\; \sum_w p(w)\,\bigl(-\log q(w)\bigr) \;=\; -\sum_w p(w)\,\log q(w).
\]
Look closely at the two distributions in that one formula. The surprise $-\log q(w)$ is measured
with the \emph{model's} probabilities, but the averaging $\sum_w p(w)\,(\cdots)$ is done with
\emph{reality's} frequencies. That is the one difference from entropy, where \emph{both} were the
true $p$. This mixed quantity --- real-world averaging of model-felt surprise --- has a name, the
\textbf{cross-entropy}, but the name matters less than the picture: \emph{training lowers the
model's average surprise at real language.}

\section{Two parts, and only one of them can move}

That average surprise is secretly \emph{two things added together}.

\textbf{Part one: unavoidable surprise.} Some of it is nobody's fault. After ``I flipped a fair
coin and it came up \underline{\hspace{1.5em}},'' even a perfect predictor can only say 50/50 ---
it \emph{will} be surprised one bit's worth every time, and no cleverness removes that. Language
is full of this: after ``I went to the \underline{\hspace{1.5em}},'' many words are genuinely
possible. This floor is exactly the \textbf{entropy} $H(p)$ of language from the last section:
the average surprise that remains even with perfect probabilities.

\textbf{Part two: avoidable surprise.} The extra surprise the model suffers \emph{only because
its guesses are off} --- not because language is hard, but because the model is wrong. A worse
model is surprised more than it needs to be. This is the part a better model shrinks.

\[
  \text{average surprise} \;=\; \underbrace{\text{entropy of language}}_{\text{fixed}}
    \;+\; \underbrace{\text{the model's avoidable mistakes}}_{\text{all that can change}}.
\]

That second part --- the gap between the model and reality --- is the \textbf{KL-divergence}.
Stripped of the name it is just ``how much extra surprise the model eats purely for being
wrong.'' It is zero only when the model's bets exactly match the true odds of language, and a
positive number otherwise; the more wrong, the bigger it is.

\textbf{The punchline.} The entropy of language does not depend on the model at all --- it is a
fixed number, a property of the data. So when training drives the total down, the fixed part
cannot budge, and \emph{the only thing that can actually move is the gap}. Lowering the loss and
shrinking the gap are therefore the same act.
\begin{summary}
Training a language model is dragging the model's picture of language closer to the real thing,
and the KL-divergence is the tape measure for how far apart they still are.
\end{summary}

\section{The two formulas look nothing alike --- here is the bridge}

So far, so verbal. But written down, the two quantities do not look like each other, and it is
worth seeing exactly why they are secretly the same.

The \textbf{loss} is the cross-entropy built in the last section but one:
\[
  \text{loss} \;=\; -\sum_w p(w)\,\log q(w),
\]
a single bare logarithm, mixing the true frequencies $p$ and the model's guesses $q$.

The \textbf{KL-divergence} is written with a \emph{ratio} inside the logarithm, which is what
makes it look like a different animal:
\[
  D_{\mathrm{KL}}(p \,\|\, q) \;=\; \sum_w p(w)\,\log\frac{p(w)}{q(w)}.
\]
A fraction $p/q$ where the loss had a bare $q$. They genuinely do not resemble one another --- but
they are one step apart, and the step is the very property of the logarithm that gave us surprise
in the first place. There it turned a product into a sum; here it turns a quotient into a
difference:
\[
  \log\frac{p}{q} \;=\; \log p - \log q.
\]
Substitute that into the KL formula and the single sum splits into two:
\[
  D_{\mathrm{KL}}(p \,\|\, q)
    \;=\; \sum_w p(w)\,\bigl(\log p(w) - \log q(w)\bigr)
    \;=\; \underbrace{\sum_w p(w)\,\log p(w)}_{(\star)}
      \;\underbrace{-\;\sum_w p(w)\,\log q(w)}_{(\star\star)}.
\]
Now read the two halves against the earlier sections.
\begin{itemize}[leftmargin=1.4em]
  \item The second half $(\star\star)$, $-\sum_w p(w)\log q(w)$, is \emph{exactly the loss} ---
  letter for letter. It is the average surprise.
  \item The first half $(\star)$, $\sum_w p(w)\log p(w)$, uses only the true frequencies $p$ and
  never mentions the model $q$. It is $-H(p)$: \emph{minus the entropy} of language --- the fixed
  floor of Section~2. Nothing the model does can change it.
\end{itemize}
So the equation reads, in the words of the last section,
\[
  \underbrace{D_{\mathrm{KL}}(p\|q)}_{\text{the gap}}
    \;=\; \underbrace{-\,H(p)}_{\text{a fixed constant}}
      \;+\; \underbrace{\text{loss}}_{\text{average surprise}},
  \qquad\text{i.e.}\qquad
  \text{loss} \;=\; H(p) \;+\; D_{\mathrm{KL}}(p\|q).
\]
The loss formula and the KL formula are not two measurements. The loss is \emph{literally the
second half of the KL}. The half that looked unfamiliar --- the ratio, the extra $\log p$ ---
collects into one fixed number, the entropy, that the model cannot move. That the two differ by
nothing but this fixed amount is the whole reason training the one is training the other --- which
the next section makes exact.

\begin{insight}
There is a small sleight-of-hand worth savouring. To \emph{measure} the gap you would seem to
need $p$ --- the true odds of language, which nobody knows. But the only place $p$ appears on its
own is the fixed floor $(\star) = -H(p)$, and a fixed number has no effect on which way is
downhill. The part that actually guides training is $(\star\star)$, the average of $-\log q$,
which needs only the model's own probabilities for words that really occurred --- and those you
can read off directly. You descend a gap you could never compute, because its uncomputable part
never moves.
\end{insight}

\begin{insight}
One asymmetry to file away. The KL above is written $D_{\mathrm{KL}}(p\|q)$, with the \emph{true}
distribution $p$ in front --- the average is over words that genuinely occur. Written the other
way round, $D_{\mathrm{KL}}(q\|p)$, it is a different number; KL is a directed gap, not a
symmetric distance. For ordinary training it is always the first form, because the examples are
drawn from real language, so the averaging is over $p$ whether we name it or not.
\end{insight}

\section{Minimising the loss is minimising the KL}

Write the model's distribution as $q_\theta$, to keep in view that it depends on the model's
tunable parameters $\theta$ --- its weights, the only things training is allowed to change. The
decomposition from the last section becomes
\[
  \text{loss}(\theta) \;=\; H(p) \;+\; D_{\mathrm{KL}}(p \,\|\, q_\theta).
\]
Everything now turns on a single observation: \textbf{$H(p)$ does not contain $\theta$.} The
entropy of language is built from the true frequencies $p$ alone; the model never appears in it.
So as training tunes $\theta$, the term $H(p)$ sits perfectly still --- a constant offset --- and
every bit of the $\theta$-dependence lives in the KL term.

\textbf{Same minimiser.} Two functions that differ by a constant have the same shape. Picture the
loss as a landscape over the space of all possible models: $D_{\mathrm{KL}}$ is the \emph{same}
landscape, shifted bodily upward by $H(p)$. Raising or lowering a landscape as a whole moves none
of its features --- the valleys stay where they are, and in particular the lowest point stays
where it is. So the $\theta$ that minimises the loss is exactly the $\theta$ that minimises the
KL.

\textbf{Same route.} It is not only the destination that agrees; so does every step toward it.
Training moves downhill along the gradient, and the gradient of a constant is zero, so
\[
  \frac{\partial\,\text{loss}}{\partial \theta}
    \;=\; \underbrace{\frac{\partial\, H(p)}{\partial\theta}}_{=\,0}
      \;+\; \frac{\partial\, D_{\mathrm{KL}}(p\|q_\theta)}{\partial\theta}
    \;=\; \frac{\partial\, D_{\mathrm{KL}}(p\|q_\theta)}{\partial\theta}.
\]
At every point in the landscape the two have the identical slope. Each training step is the same
step whether you call it ``descend the loss'' or ``descend the KL'' --- not merely the same
endpoint, but the same path, nudge for nudge.

\begin{summary}
Minimising the loss \emph{is} minimising the KL-divergence --- the same minimiser, and the same
gradient at every step --- because the two differ only by $H(p)$, the entropy of language, which
the model cannot touch.
\end{summary}

This also says, sharply, what training is reaching for. The KL is never negative, and it is zero
in exactly one situation: when $q_\theta = p$, the model's distribution equal to language's own.
So the best the model can possibly do is drive the KL to $0$ --- become a perfect mirror of
language --- and at that point the loss is not $0$ but
\[
  \text{loss} \;=\; H(p) + 0 \;=\; H(p).
\]
The entropy of language is a \textbf{floor} the loss can approach but never cross. A model is not
trying to push its loss to zero --- impossible, unless language itself were perfectly predictable
--- it is trying to push the \emph{gap above the floor} to zero. The surprise left over at the
end, $H(p)$, was never the model's to remove: it is the irreducible uncertainty of language, and
reaching it means the model has learned everything there was to learn.

\begin{insight}
Concretely, take the biased $0.9$ coin of Section~2, whose entropy is $0.47$ bits. A model
learning it cannot drive its loss below $0.47$ bits, however perfectly it trains --- that number
is the coin's own randomness, not the model's shortfall. A loss of $0.47$ bits does not mean the
model is still $0.47$ short of the answer; it means it has the answer \emph{exactly}, $q = p$, and
the residue is the coin. So when you read a training curve, the number that matters is not how
close the loss is to $0$ but how close it is to $H(p)$: the distance down to that floor is the KL
still left to remove.
\end{insight}

\end{document}
